% =============================================================================
\chapter{Admin}
\label{ch:admin}
% =============================================================================

Seven prominent admin pages: Pipelines, Infrastructure, Cloud GPUs, Storage,
Naming Rules, Queue Manager and Output Profiles. We drill into
\textbf{x121} for pipeline detail, \textbf{RunPod} for cloud GPU detail, and
\textbf{delivery\_video} for the naming-rule edit modal. Other
\texttt{/admin/*} routes are out of scope.

% --------------------------------------------------------------------------
\section{Pipelines}
\label{sec:admin-pipelines}
% Route: /admin/pipelines
% Component: features/admin/pipelines/PipelineListPage.tsx

Pipelines are first-class top-level entities: a pipeline captures its own seed
image requirements, workflows, tracks, naming rules and delivery formats. New
content verticals are added by creating a pipeline record (seed slots, workflows,
naming) rather than by forking the codebase.

\screenshot{08-admin-pipelines-list-hero.png}{Pipelines list — one card per pipeline.}

\begin{itemize}
  \pagelement{Pipeline cards}{Code, name, owner, state, quick-edit link.}
  \pagelement{Create action}{Opens the pipeline creation flow (UI shown, not submitted).}
\end{itemize}

\screenshot{08-admin-pipelines-x121-detail.png}{\texttt{x121} pipeline detail / edit page.}

\begin{itemize}
  \pagelement{General}{Code, name, owner, description.}
  \pagelement{Seed slots}{Per-slot config (name, required flag).}
  \pagelement{Hooks}{Pipeline hook registrations.}
  \pagelement{Save / Delete}{Captured but never submitted.}
\end{itemize}

% --------------------------------------------------------------------------
\section{Infrastructure}
\label{sec:admin-infrastructure}
% Route: /admin/infrastructure
% Component: features/admin/infrastructure/InfrastructureControlPanel.tsx

Single operational command centre for cloud GPU and ComfyUI infrastructure. It
replaces two older UIs (an inline panel on the Generation page and the
read-only cloud-GPU dashboard) with one page that covers the full lifecycle:
provision $\rightarrow$ SSH startup $\rightarrow$ ComfyUI readiness $\rightarrow$
WebSocket connection $\rightarrow$ queue distribution $\rightarrow$ teardown.
Includes orphan detection and failed-connection recovery.

\screenshot{08-admin-infrastructure-hero.png}{Infrastructure control panel.}

\begin{itemize}
  \pagelement{Provider management}{Summary of registered cloud providers.}
  \pagelement{Provision wizard}{Opens the provisioning flow (UI shown, not submitted).}
  \pagelement{Instance actions}{Start / stop / reclaim buttons per instance.}
  \pagelement{Orphan panel}{Orphaned instances to reclaim.}
  \pagelement{Activity log}{Recent provisioning events.}
\end{itemize}

% --------------------------------------------------------------------------
\section{Cloud GPUs}
\label{sec:admin-cloud-gpus}
% Route: /admin/cloud-gpus
% Component: features/admin/cloud-gpus/CloudGpuDashboard.tsx

Provider registry and per-provider configuration. The integration layer is
built around a provider trait (RunPod Pods and RunPod Serverless are the first
implementations) so Vast.ai, Lambda Labs, etc. can be added without refactoring
core queueing or dispatch.

\screenshot{08-admin-cloud-gpus-hero.png}{Cloud GPU dashboard.}

\begin{itemize}
  \pagelement{Stats strip}{Total providers, active providers, instance counts.}
  \pagelement{Provider list}{Card per provider with health + cost summary.}
\end{itemize}

\subsection{RunPod — provider detail}
% Route: /admin/cloud-gpus/$providerId (RunPod → id 1)
% Component: features/admin/cloud-gpus/components/CloudProviderDetail.tsx

Per-provider detail: instance list, supported GPU types, scaling rules, and
cost tracking. RunPod is the reference provider implementation.

\screenshot{08-admin-cloud-gpus-runpod-hero.png}{RunPod provider detail — default (Instances) tab.}

\begin{itemize}
  \pagelement{Four tabs}{Instances, GPU Types, Scaling Rules, Cost.}
\end{itemize}

\subsubsection{Instances tab}
\screenshot{08-admin-cloud-gpus-runpod-instances-tab.png}{RunPod / Instances tab.}

\subsubsection{GPU Types tab}
\screenshot{08-admin-cloud-gpus-runpod-gpu-types-tab.png}{RunPod / GPU Types tab.}

\subsubsection{Scaling Rules tab}
\screenshot{08-admin-cloud-gpus-runpod-scaling-rules-tab.png}{RunPod / Scaling Rules tab.}

\subsubsection{Cost tab}
\screenshot{08-admin-cloud-gpus-runpod-cost-tab.png}{RunPod / Cost tab.}

% --------------------------------------------------------------------------
\section{Storage}
\label{sec:admin-storage}
% Route: /admin/storage
% Component: features/storage/StoragePage.tsx

Backend configuration for pluggable storage. File I/O goes through a
\texttt{StorageProvider} trait with local-filesystem and S3-compatible
implementations (MinIO, DO Spaces, Backblaze B2), with policy-driven hot / cold
tiering between them. Each pipeline's imports land in its own isolated storage
path via the pipeline-scoped key helper.

\screenshot{08-admin-storage-hero.png}{Storage backends list.}

\begin{itemize}
  \pagelement{Backend table}{Current backends, kind (local / S3 / Azure), status.}
  \pagelement{Add backend}{Opens \texttt{BackendFormModal}. UI shown, never submitted.}
\end{itemize}

\screenshotnarrow{08-admin-storage-backend-form-modal.png}{BackendFormModal — new-backend form.}

% --------------------------------------------------------------------------
\section{Naming Rules}
\label{sec:admin-naming}
% Route: /admin/naming
% Component: features/naming-rules/NamingRulesPage.tsx

Centralised, token-based naming engine. Every file the platform produces (scene
videos, variants, thumbnails, delivery ZIPs, exports) is named from templates
stored per category in the database, with live preview. Global by default with
per-project overrides for client-specific delivery specs — replacing 7+
hardcoded patterns that previously required code changes.

\screenshot{08-admin-naming-hero.png}{Admin Naming Rules — category grid.}

\begin{itemize}
  \pagelement{Category groups}{Categories grouped by domain (Generation, Storage, Export, Delivery).}
  \pagelement{Category card}{Shows current template and an example output.}
  \pagelement{Edit action}{Clicking opens a modal with the template editor (captured below).}
\end{itemize}

\screenshot{08-admin-naming-delivery-video-modal.png}{Edit modal for the \texttt{delivery\_video} category — template input, token chips, live preview.}

\begin{itemize}
  \pagelement{Template pattern}{Textarea for the naming template.}
  \pagelement{Token chips}{Clickable token chips insert at cursor.}
  \pagelement{Preview}{Live preview of a rendered filename.}
  \pagelement{Save / Cancel}{Captured but never saved.}
\end{itemize}

% --------------------------------------------------------------------------
\section{Queue Manager}
\label{sec:admin-queue-manager}
% Route: /admin/queue
% Component: features/queue/QueueManagerPage.tsx

Admin-only operational view over the same job queue the production page shows,
with added controls: reassign jobs across workers, drain an instance without
killing in-flight jobs, move jobs to the front, and redistribute load. This is
what turns the queue from a status page into something admins can actively
manage.

\screenshot{08-admin-queue-manager-hero.png}{Admin Queue Manager — queue state + activity.}

\begin{itemize}
  \pagelement{Queue status view}{All jobs across the system with admin-only controls.}
  \pagelement{Activity log}{Global queue event stream.}
\end{itemize}

% --------------------------------------------------------------------------
\section{Output Profiles}
\label{sec:admin-output-profiles}
% Route: /admin/output-profiles
% Component: features/output-profiles/OutputProfilesPage.tsx

Platform-wide delivery encoding presets (resolution, codec, container, bitrate,
framerate). One profile is designated as the platform default; projects can
override it with a project-specific default so the ExportPanel auto-selects
the right profile without asking.

\screenshot{08-admin-output-profiles-hero.png}{Output Profiles list.}

\begin{itemize}
  \pagelement{Profile cards}{Each profile: name, codec, container, bitrate, resolution.}
  \pagelement{Default profile}{Captured below in detail.}
  \pagelement{Create / Edit}{Form opens in a side panel (captured, never submitted).}
\end{itemize}

\screenshotnarrow{08-admin-output-profiles-default-detail.png}{Default output profile — edit panel.}
